\chapter{Introdução}

O estudo da resposta de circuitos lineares em regime permanente senoidal é fundamental para a compreensão do comportamento de sistemas no domínio da frequência \cite{nilsson}. 

Nesta prática de laboratório, realizou-se a análise teórica e experimental da influência da função de transferência na resposta de um filtro ativo de realimentação múltipla (MFB) implementado com amplificador operacional. O trabalho foi estruturado em etapas complementares: a análise analítica e simbólica utilizando o software Maxima \cite{maxima}, a simulação computacional para validação dos modelos via LTSpice \cite{ltspice}, e as medições experimentais em bancada utilizando osciloscópio para caracterização de magnitude e fase em regime permanente.

\section{Objetivos}

\begin{itemize}
	\item Exercitar o cálculo analítico da função de transferência $H(s)$ para circuitos lineares complexos[cite: 22].
	\item Aplicar a função de transferência para determinar a resposta em regime permanente senoidal (RPS) de um circuito linear submetido a entradas senoidais.
	\item Praticar o uso de softwares de álgebra simbólica para a resolução e simplificação de equações de circuitos elétricos \cite{maxima}.
	\item Utilizar ferramentas de simulação para verificação prévia de resultados teóricos \cite{ltspice}.
	\item Desenvolver competências práticas no uso do osciloscópio para medição de amplitudes e defasagem entre sinais de entrada e saída.
\end{itemize}

O procedimento consistiu na modelagem matemática do circuito, seguida pela coleta de dados experimentais em uma faixa de frequências de 40 Hz a 11 kHz, visando comparar o comportamento do filtro real com as predições teóricas obtidas.